%% sol-ch05.tex — Solutions for Chapter 5: Closure Properties

\chapter{Closure Properties --- Solutions}

\begin{enumerate}[{5.}1. ]

%% 5.1
\item Let $F_\Sigma$ be the collection of all finite languages over $\Sigma$.

\begin{enumerate}[(a)]
  \item \textbf{Complementation: No.}  $\Sigma^*$ is infinite, so for any finite $L$,
  $\Sigma^*\setminus L$ is infinite (assuming $|\Sigma|\ge 1$), hence not in $F_\Sigma$.

  \item \textbf{Union: Yes.}  If $L_1,L_2$ are finite, then $|L_1\cup L_2|\le|L_1|+|L_2|<\infty$.

  \item \textbf{Intersection: Yes.}  $|L_1\cap L_2|\le|L_1|<\infty$.

  \item \textbf{Concatenation: Yes.}  $|L_1\cdot L_2|\le|L_1|\cdot|L_2|<\infty$.

  \item \textbf{Kleene closure: No.}  Let $L=\{\pmb{a}\}$.  Then $L^*=\{\pmb{a}\}^*=\{\lambda,\pmb{a},\pmb{aa},\ldots\}$ is infinite.

  \item \textbf{Relative complement: Yes.}  $L_1\setminus L_2\subseteq L_1$, so $|L_1\setminus L_2|\le|L_1|<\infty$.
\end{enumerate}

%% 5.2
\item Let $C_\Sigma$ be the collection of all cofinite languages (complements of finite languages).

\begin{enumerate}[(a)]
  \item \textbf{Complementation: No.}  Let $L=\Sigma^*\setminus\{\pmb{a}\}$.  Then
  $L\in C_\Sigma$, but $\sim L=\{\pmb{a}\}$ is finite, hence not in $C_\Sigma$.

  \item \textbf{Union: Yes.}  If $L_1=\Sigma^*\setminus F_1$ and $L_2=\Sigma^*\setminus F_2$, then $L_1\cup L_2=\Sigma^*\setminus(F_1\cap F_2)$.  Since $F_1\cap F_2$ is finite, $L_1\cup L_2\in C_\Sigma$.

  \item \textbf{Intersection: Yes.}  $L_1\cap L_2=\Sigma^*\setminus(F_1\cup F_2)$.  Since $F_1$ and $F_2$ are finite, so is $F_1\cup F_2$, hence $L_1\cap L_2\in C_\Sigma$.

  \item \textbf{Concatenation: Yes.}  If $L_1,L_2\in C_\Sigma$ (both infinite), then $L_1\cdot L_2$ contains arbitrarily long words, but let us verify cofiniteness.  Since $\Sigma^*\setminus(L_1\cdot L_2)\subseteq$ the set of words not writable as a product from $L_1$ and $L_2$; this set is finite because $L_1$ and $L_2$ each omit only finitely many words.  More precisely: there exists $N_1$ such that all words of length $\ge N_1$ are in $L_1$, and similarly $N_2$ for $L_2$.  Then any word of length $\ge N_1+N_2$ can be split so that both parts are in $L_1$ and $L_2$, so $L_1\cdot L_2$ is cofinite.

  \item \textbf{Kleene closure: Yes.}  Since $L$ is cofinite, there is an $N$
  such that every word of length at least $N$ lies in $L$.  Hence every word of
  length at least $N$ also lies in $L^*$ (just take it as a one-factor product).
  Therefore $\Sigma^*\setminus L^*$ contains only words of length $<N$, so it is
  finite and $L^*\in C_\Sigma$.

  \item \textbf{Relative complement: No.}  $L_1\setminus L_2=L_1\cap\sim L_2$.  Here $\sim L_2$ is finite, so $L_1\setminus L_2\subseteq\sim L_2$ is finite, not cofinite in general.  Counterexample: $L_1=\Sigma^*\setminus\emptyset=\Sigma^*$ and $L_2=\Sigma^*\setminus\{\pmb{a}\}$.  Then $L_1\setminus L_2=\{\pmb{a}\}$, which is finite, not cofinite.
\end{enumerate}

%% 5.3
\item $B_\Sigma=F_\Sigma\cup C_\Sigma$.

\begin{enumerate}[(a)]
  \item \textbf{Complementation: Yes.}  The complement of a finite language is cofinite and vice versa.
  \item \textbf{Union: Yes.}  If both $L_1,L_2$ are finite, $L_1\cup L_2$ is finite.  If one is cofinite, $L_1\cup L_2$ is cofinite.  If both are cofinite, $L_1\cup L_2$ is cofinite (by 5.2b).  \textbf{Closed.}
  \item \textbf{Intersection: Yes.}  If both finite: finite.  If one cofinite and one finite: finite.  If both cofinite: cofinite (by 5.2c).
  \item \textbf{Concatenation: It depends on $\Sigma$.}
  If $|\Sigma|=1$, say $\Sigma=\{\pmb{a}\}$, then $B_\Sigma$ \emph{is} closed
  under concatenation.  Finite unary languages correspond to finite sets of
  lengths, and cofinite unary languages correspond to cofinite sets of lengths.
  The sum of two finite sets of lengths is finite; the sum of a nonempty finite
  set and a cofinite set is cofinite; and the sum of two cofinite sets is
  cofinite.  So every concatenation stays in $B_\Sigma$.

  If $|\Sigma|\ge 2$, then $B_\Sigma$ is \emph{not} closed under concatenation.
  For example, let $L_1=\{\pmb{a}\}$ and
  $L_2=\Sigma^*\setminus\{\pmb{aa}\}$.  Then $L_1\in F_\Sigma$ and
  $L_2\in C_\Sigma$, but
  \[
  L_1\cdot L_2=\{\pmb{a}w\mid w\neq \pmb{aa}\},
  \]
  the set of all words beginning with $\pmb{a}$ except $\pmb{aaa}$.  This
  language is infinite, and its complement is also infinite, so it is neither
  finite nor cofinite.

  Therefore the correct answer is:
  \[
  |\Sigma|=1 \Rightarrow \text{closed},\qquad |\Sigma|\ge 2 \Rightarrow \text{not closed}.
  \]
  \item \textbf{Kleene closure: No.}  Pick any symbol $\pmb{a}\in\Sigma$ and let
  $L=\{\pmb{aa}\}$.  Then $L\in F_\Sigma\subseteq B_\Sigma$, but
  \[
  L^*=\{(\pmb{aa})^n\mid n\ge 0\}.
  \]
  This language is infinite.  If $|\Sigma|=1$, its complement is the set of
  odd-length words, which is infinite.  If $|\Sigma|\ge 2$, its complement
  contains every word with some symbol other than $\pmb{a}$, so it is again
  infinite.  Thus $L^*$ is neither finite nor cofinite, so $B_\Sigma$ is not
  closed under Kleene closure.
  \item \textbf{Relative complement: Yes.}  If $L_1\in C_\Sigma$ and
  $L_2\in C_\Sigma$, then $L_1\setminus L_2=L_1\cap\sim L_2$ where
  $\sim L_2$ is finite, so the result is finite and therefore belongs to
  $F_\Sigma\subseteq B_\Sigma$.  If $L_1\in F_\Sigma$, then
  $L_1\setminus L_2\subseteq L_1$ is finite.  Hence $B_\Sigma$ is closed under
  relative complement.
\end{enumerate}

%% 5.4
\item $I_\Sigma=\wp(\Sigma^*)\setminus F_\Sigma$, the infinite languages.

\begin{enumerate}[(a)]
  \item \textbf{Complementation: No.}  $\Sigma^*\setminus L$ may be finite.  E.g., $L=\Sigma^*\setminus\{\lambda\}$ is infinite; $\sim L=\{\lambda\}$ is finite.
  \item \textbf{Union: Yes.}  Union of two infinite sets is infinite.
  \item \textbf{Intersection: No.}  $L_1=\{\pmb{a}^{2n}\mid n\ge 0\}$ and $L_2=\{\pmb{a}^{2n+1}\mid n\ge 0\}$ are both infinite, but $L_1\cap L_2=\emptyset$, which is finite.
  \item \textbf{Concatenation: Yes.}  If $L_1,L_2$ are infinite, pick any fixed
  word $y\in L_2$.  Then
  \[
  \{xy\mid x\in L_1\}\subseteq L_1\cdot L_2.
  \]
  The map $x\mapsto xy$ is injective, so $\{xy\mid x\in L_1\}$ is infinite.
  Therefore $L_1\cdot L_2$ is infinite.
  \item \textbf{Kleene closure: Yes.}  $L^*\supseteq L$ when $\lambda\in L^*$, and $L^*$ is always infinite if $L$ contains any nonempty word.
  \item \textbf{Relative complement: No.}  $L_1\setminus L_2$ could be finite even if $L_1,L_2$ are infinite (e.g., $L_1=L_2$).
\end{enumerate}

%% 5.5
\item $J_\Sigma=\wp(\Sigma^*)\setminus C_\Sigma$, languages with infinite complements.

\begin{enumerate}[(a)]
  \item \textbf{Complementation: No.}  $L\in J_\Sigma$ means $\sim L$ is infinite, i.e., $\sim L\in I_\Sigma$.  Then $\sim(\sim L)=L$; is $\sim L\in J_\Sigma$?  $\sim(\sim L)=L$ is infinite (since $L\in J_\Sigma$ is not cofinite, its complement is infinite, so $L$ itself need not be infinite).  Counterexample: $L=\emptyset\in J_\Sigma$ (complement is $\Sigma^*$, infinite), but $\sim L=\Sigma^*\not in J_\Sigma$ (complement is $\emptyset$, finite).  \textbf{Not closed.}
  \item \textbf{Union: No.}  Over $\Sigma=\{\pmb{a}\}$, let
  $L_1=\{\pmb{a}^{2n}\mid n\ge 0\}$ and $L_2=\{\pmb{a}^{2n+1}\mid n\ge 0\}$.
  Both have infinite complements, so both belong to $J_\Sigma$.  But
  $L_1\cup L_2=\Sigma^*$, whose complement is finite, so $L_1\cup L_2\not in J_\Sigma$.
  \item \textbf{Intersection: Yes.}  If $\sim L_1$ and $\sim L_2$ are both infinite, then $\sim(L_1\cap L_2)=\sim L_1\cup\sim L_2\supseteq\sim L_1$, which is infinite.
  \item \textbf{Concatenation: No.}  Over $\Sigma=\{\pmb{a}\}$, let $L_1=\{\pmb{a}^n\mid n\text{ odd or }n=2\}$ and $L_2=\{\pmb{a}^{2k}\mid k\ge 1\}$.  Then $\sim L_1=\{\lambda,\pmb{a}^4,\pmb{a}^6,\ldots\}$ and $\sim L_2=\{\lambda,\pmb{a},\pmb{a}^3,\pmb{a}^5,\ldots\}$ are both infinite, so $L_1,L_2\in J_\Sigma$.  But $L_1\cdot L_2=\{\pmb{a}^n\mid n\ge 3\}$ (any $n\ge 3$ is expressible as odd$+$even or $2+$even), whose complement $\{\lambda,\pmb{a},\pmb{a}^2\}$ is finite---so $L_1\cdot L_2\not in J_\Sigma$.
  \item \textbf{Kleene closure: No.}  Over $\Sigma=\{\pmb{a}\}$, let $L=\{\pmb{a}^n\mid n\text{ odd}\}$.  Then $\sim L=\{\lambda,\pmb{a}^2,\pmb{a}^4,\ldots\}$ is infinite, so $L\in J_\Sigma$.  But odd$+$odd$=$even, even$+$odd$=$odd, so $L^k$ contains all $\pmb{a}$-strings of the right parity at length $\ge k$; inductively $L^*=\{\pmb{a}\}^*=\Sigma^*$.  Then $\sim L^*=\emptyset$ is finite, so $L^*\not in J_\Sigma$.
  \item \textbf{Relative complement: Yes.}  $\sim(L_1\setminus L_2)=\sim(L_1\cap\sim L_2)=\sim L_1\cup L_2\supseteq\sim L_1$.  Since $L_1\in J_\Sigma$, $\sim L_1$ is infinite, so $\sim(L_1\setminus L_2)$ is infinite, hence $L_1\setminus L_2\in J_\Sigma$. \qed
\end{enumerate}

%% 5.6
\item $\pmb{E}$: all languages over $\{\pmb{a},\pmb{b}\}$ that contain $\pmb{abba}$.

\begin{enumerate}[(a)]
  \item \textbf{Complementation: No.}  $\sim L$ contains all words \emph{not} in $L$; since $L$ contains $\pmb{abba}$, $\sim L$ does not contain $\pmb{abba}$, so $\sim L\not in\pmb{E}$.
  \item \textbf{Union: Yes.}  If $\pmb{abba}\in L_1$, then $\pmb{abba}\in L_1\cup L_2$.
  \item \textbf{Intersection: Yes.}  If $L_1,L_2\in\pmb{E}$ then both contain $\pmb{abba}$, so $L_1\cap L_2$ also contains $\pmb{abba}$.  \textbf{Closed.}
  \item \textbf{Concatenation: No.}  Here $\pmb{E}$ means languages that have
  the specific word $\pmb{abba}$ as an element.  Let
  $L_1=L_2=\{\pmb{abba}\}$.  Then $L_1,L_2\in\pmb{E}$, but
  $L_1\cdot L_2=\{\pmb{abbaabba}\}$, which does not contain $\pmb{abba}$ as a
  member.  So $\pmb{E}$ is not closed under concatenation.
  \item \textbf{Kleene closure: Yes.}  $L^*$ contains $L^1=L$, so $\pmb{abba}\in L\subseteq L^*$.
  \item \textbf{Relative complement: No.}  $L_1\setminus L_2$ might not contain $\pmb{abba}$ if $\pmb{abba}\in L_2$.  E.g., $L_1=L_2=\{\pmb{abba}\}$: $L_1\setminus L_2=\emptyset\not in\pmb{E}$.
\end{enumerate}

%% 5.7
\item \textbf{Intersection closed does not imply union closed.}

\textbf{Counterexample.}  Let $\mathcal{C}=\{\emptyset\}\cup\{\{x\}\mid x\in\Sigma^*\}$ (the empty language together with all singleton languages).

\emph{Closed under intersection}: $\{x\}\cap\{y\}=\{x\}$ if $x=y$, and $\emptyset$ if $x\neq y$; both are in $\mathcal{C}$.

\emph{Not closed under union}: $\{x\}\cup\{y\}=\{x,y\}$ for $x\neq y$, which is a two-element set and not in $\mathcal{C}$. $\square$

%% 5.8
\item \textbf{If closed under intersection and complement, then closed under union (by De Morgan).}

If $\mathcal{C}$ is closed under intersection and complement, then for any $L_1,L_2\in\mathcal{C}$:
$\sim L_1,\sim L_2\in\mathcal{C}$ (by closure under complement), so
$\sim L_1\cap\sim L_2\in\mathcal{C}$ (by closure under intersection), so
$\sim(\sim L_1\cap\sim L_2)=L_1\cup L_2\in\mathcal{C}$ (by De Morgan and closure under complement).  Hence $\mathcal{C}$ is also closed under union. \qed

%% 5.9
\item \textbf{Closed under concatenation does not imply closed under Kleene closure.}

Let $\mathcal{C}=F_\Sigma$ (all finite languages).  Closed under concatenation (5.1d).  Not closed under Kleene closure: $\{\pmb{a}\}^*$ is infinite (5.1e).

%% 5.10
\item \textbf{Closed under Kleene closure does not imply closed under concatenation.}

Let
\[
\mathcal{C}=\{\{w\}^*\mid w\in\Sigma^*\}.
\]
This includes $\{\lambda\}=\{\lambda\}^*$ (take $w=\lambda$).  Each language in
$\mathcal{C}$ is closed under Kleene closure, because
\[
(\{w\}^*)^*=\{w\}^*.
\]
But $\mathcal{C}$ is not closed under concatenation: over
$\Sigma=\{\pmb{a},\pmb{b}\}$,
\[
\{\pmb{a}\}^*\cdot\{\pmb{b}\}^*=\{\pmb{a}^m\pmb{b}^n\mid m,n\ge 0\},
\]
and this language is not of the form $\{u\}^*$ for any word $u$.

%% 5.11
\item \textbf{Closed under complement does not imply closed under relative complement.}

\textbf{Counterexample.}  Let
\[
K=\{\emptyset,\Sigma^*\}\cup\{\{w\}\mid w\in\Sigma^*\}\cup
\{\Sigma^*\setminus\{w\}\mid w\in\Sigma^*\}.
\]
This collection is closed under complement by construction.  But it is not
closed under relative complement: if $u\neq v$, then
\[
(\Sigma^*\setminus\{u\})\setminus\{v\}=\Sigma^*\setminus\{u,v\}\not in K.
\]
So closure under complement does not imply closure under relative complement.

%% 5.12
\item \textbf{A finite set of numbers closed under $\sqrt{\ }$:}

$\{0,1\}$.  $\sqrt{0}=0\in\{0,1\}$ and $\sqrt{1}=1\in\{0,1\}$.

%% 5.13
\item \textbf{An infinite set of numbers closed under $\sqrt{\ }$:}

$[0,1]=\{x\in\mathbb{R}\mid 0\le x\le 1\}$.  For any $x\in[0,1]$, $\sqrt{x}\in[0,1]$ since $0\le x\le 1\Rightarrow 0\le\sqrt{x}\le 1$.

%% 5.14
\item \textbf{Algorithm for deterministic union machine $A^U$.}

Given $A_1=\langle\Sigma,S_1,s^1_0,\delta_1,F_1\rangle$ and $A_2=\langle\Sigma,S_2,s^2_0,\delta_2,F_2\rangle$ (both DFAs, $S_1\cap S_2=\emptyset$):

1.  Build the NDFA $A^\cup=\langle\Sigma,S_1\cup S_2,\{s^1_0,s^2_0\},\delta^\cup,F_1\cup F_2\rangle$
    (as in Theorem~5.2).
2.  Apply Definition~4.5 (subset construction) to $A^\cup$ to produce a DFA $A^U$.
    Start state: $\{s^1_0,s^2_0\}$.
    For each reachable state-set $Q\subseteq S_1\cup S_2$ and each $\pmb{a}\in\Sigma$:
    $\delta^U(Q,\pmb{a})=\delta^d(Q,\pmb{a})=\bigcup_{q\in Q}\delta^\cup(q,\pmb{a})$.
    Final states: all $Q$ with $Q\cap(F_1\cup F_2)\neq\emptyset$.

%% 5.15
\item \begin{enumerate}[(a)]
  \item \textbf{Direct product machine for union.}  Define
  $A^u=\langle\Sigma,S_1\times S_2,\langle s^1_0,s^2_0\rangle,\delta^u,F^u\rangle$ where
  $\delta^u(\langle s,t\rangle,\pmb{a})=\langle\delta_1(s,\pmb{a}),\delta_2(t,\pmb{a})\rangle$
  and $F^u=\{\langle s,t\rangle\mid s\in F_1\vee t\in F_2\}$.

  \item \textbf{Proof.}  By induction, $\overline{\delta^u}(\langle s^1_0,s^2_0\rangle,x)
  =\langle\deltabar{1}(s^1_0,x),\deltabar{2}(s^2_0,x)\rangle$.
  Then $x\in L(A^u)$ iff $\overline{\delta^u}(\langle s^1_0,s^2_0\rangle,x)\in F^u$
  iff $\deltabar{1}(s^1_0,x)\in F_1$ or $\deltabar{2}(s^2_0,x)\in F_2$
  iff $x\in L(A_1)\cup L(A_2)$.

  \item \textbf{State counts (no minimization).}
  $A^\cup$ (NDFA union): $|S_1|+|S_2|$ states (plus the new start state if added separately, but in Theorem~5.2 no new state is needed).
  $A^U$ (subset construction of $A^\cup$): up to $2^{|S_1|+|S_2|}$ states, but in practice far fewer.
  $A^u$ (product): $|S_1|\cdot|S_2|=nm$ states.
  With $n=|S_1|$ and $m=|S_2|$: $A^\cup$ has $n+m$ states; $A^U$ has up to $2^{n+m}$ states; $A^u$ has $nm$ states.
\end{enumerate}

%% 5.16
\item \textbf{\DSIG\ is closed under the prefix operator $P$.}

\begin{proof}
Let $L\in\DSIG$ and $A=\langle\Sigma,S,s_0,\delta,F\rangle$ be a DFA with $L(A)=L$.
Define $A'=\langle\Sigma,S,s_0,\delta,F'\rangle$ where
$F'=\{s\in S\mid\exists y\in\Sigma^*\ni\deltabar{}(s,y)\in F\}$.
(A state is final in $A'$ iff some continuation from it leads to acceptance in $A$.)

\emph{Claim:} $L(A')=P(L)$.

$x\in P(L)$ iff $\exists y\in\Sigma^*\ni xy\in L$
iff $\exists y\in\Sigma^*\ni\deltabar{}(s_0,xy)\in F$
iff $\exists y\in\Sigma^*\ni\deltabar{}(\deltabar{}(s_0,x),y)\in F$
iff $\deltabar{}(s_0,x)\in F'$
iff $x\in L(A')$.

Since $F'$ is computable (for each state $s$, check if any state in $F$ is reachable from $s$---a finite graph problem), $A'$ is a valid DFA and $P(L)=L(A')\in\DSIG$. \qed
\end{proof}

%% 5.17
\item \textbf{\DSIG\ is closed under the suffix operator $S$.}

\begin{enumerate}[(a)]
  \item \textbf{Construction.}  Given DFA $A=\langle\Sigma,S,s_0,\delta,F\rangle$,
  build NDFA $A'=\langle\Sigma,S,S_0',\delta,F\rangle$ where
  $S_0'=\{\deltabar{}(s_0,y)\mid y\in\Sigma^*\}$, the set of states reachable
  from $s_0$.

  \item \textbf{Proof.}  $x\in S(L)$ iff $\exists y\in\Sigma^*\ni yx\in L$
  iff $\exists y\in\Sigma^*\ni\deltabar{}(s_0,yx)\in F$
  iff $\exists s=\deltabar{}(s_0,y)\in S_0'\ni\deltabar{}(s,x)\in F$
  iff $\DBAR'(\{s\},x)\cap F\neq\emptyset$ for some $s\in S_0'$
  iff $x\in L(A')$.

  \item Since $A'$ is an NDFA over a finite state set, $L(A')\in\DSIG$ by Theorem~4.1.  Hence $S(L)\in\DSIG$. \qed
\end{enumerate}

%% 5.18
\item \textbf{\DSIG\ is closed under the center operator $C$.}

\begin{enumerate}[(a)]
  \item \textbf{Construction.}  Let
  $R=\{\deltabar{}(s_0,y)\mid y\in\Sigma^*\}$ be the set of reachable states.
  Build an NDFA $A'=\langle\Sigma,S,R,\delta,F'\rangle$ where
  $F'=\{s\in S\mid\exists y\in\Sigma^*\ni\deltabar{}(s,y)\in F\}$ (states from which $F$ is reachable).

  \item \textbf{Proof.}  $x\in C(L)$ iff $\exists y,z\ni yxz\in L$
  iff $\exists y,z\ni\deltabar{}(\deltabar{}(\deltabar{}(s_0,y),x),z)\in F$
  iff $\exists s=\deltabar{}(s_0,y)\in R$ and $\deltabar{}(s,x)\in F'$
  iff $x\in L(A')$.

  \item Since $A'$ is an NDFA, $C(L)\in\DSIG$. \qed
\end{enumerate}

%% 5.19
\item \textbf{\DSIG\ is closed under $F$ (maximal words not extended in $L$).}

Let $A=\langle\Sigma,S,s_0,\delta,F\rangle$ be a minimal DFA for $L$.
$F(L)=\{x\in L\mid \text{no nonempty word }y\text{ satisfies }xy\in L\}$.
Define $A'=\langle\Sigma,S,s_0,\delta,F''\rangle$ where
$F''=\{s\in F\mid(\forall y\in\Sigma^+)(\deltabar{}(s,y)\not in F)\}$.
Then $x\in L(A')$ iff $\deltabar{}(s_0,x)\in F''$ iff $x\in L$ and there is no
nonempty $y$ with $\deltabar{}(s_0,xy)\in F$, iff $x\in F(L)$.
So $F(L)=L(A')\in\DSIG$. \qed

%% 5.20
\item \textbf{\DSIG\ is closed under reversal.}

\begin{enumerate}[(a)]
  \item \textbf{Construction.}  Given DFA $A=\langle\Sigma,S,s_0,\delta,F\rangle$,
  build NDFA $A^r=\langle\Sigma,S,F,\delta^r,\{s_0\}\rangle$ where
  $\delta^r(s,\pmb{a})=\{t\in S\mid\delta(t,\pmb{a})=s\}$ (reverse all arrows).

  \item \textbf{Proof.}  We show $x\in L(A^r)$ iff $x^r\in L(A)$.

  By induction on $|x|$: a path $f\xrightarrow{a_1}s_1\xrightarrow{a_2}\cdots\xrightarrow{a_k}s_0$ in $A^r$ (from some $f\in F$ to $s_0$) corresponds to the path $s_0\xrightarrow{a_k}\cdots\xrightarrow{a_1}f$ in $A$.  Hence $a_1a_2\cdots a_k\in L(A^r)$ iff $a_k\cdots a_1\in L(A)$, i.e., $x\in L(A^r)$ iff $x^r\in L(A)$.

  \item $L^r=L(A^r)\in\DSIG$ (since $A^r$ is an NDFA which is equivalent to some DFA). \qed
\end{enumerate}

%% 5.21
\item \begin{enumerate}[(a)]
  \item \textbf{\DSIG\ is not closed under $G(L)=\{w^r\cdot w\mid w\in L\}$.}

  Let $L=\Sigma^*$ over $\Sigma=\{\pmb{a},\pmb{b}\}$.  Then $G(L)=\{w^rw\mid w\in\Sigma^*\}$, the set of even-length palindromes (of the form $w^rw$).  By the pumping lemma, this is not FAD: take $w=\pmb{a}^n\pmb{b}^n$; then $w^rw=\pmb{b}^n\pmb{a}^n\pmb{a}^n\pmb{b}^n\in G(L)$.  If $G(L)$ were FAD with $n$ states, pump the string $\pmb{b}^n\pmb{a}^n\pmb{a}^n\pmb{b}^n$ to get a contradiction (pumping the first block of $\pmb{b}$s destroys the palindrome structure).

  \item \textbf{\NSIG\ is closed under $G$.}  Let $L\in\NSIG$.  If $G(L)$ were
  FAD, then Theorem~5.11 would imply that $G(L)^b$ is FAD.  But
  $G(L)^b=L^r$, and Exercise~5.20 shows that reversal preserves the FAD
  languages.  Hence $L=(L^r)^r$ would be FAD, contradicting $L\in\NSIG$.
  Therefore $G(L)\in\NSIG$.

  \item \textbf{\NSIG\ is not closed under $b$.}  By Lemma~5.6,
  $L=\{\pmb{a}^n\pmb{b}^n\mid n\ge 0\}\in\NSIG$, but
  $L^b=\{\pmb{a}^n\mid n\ge 0\}$ is FAD.
\end{enumerate}

%% 5.22
\item \textbf{\NSIG\ is closed under reversal.}

Assume $L\in\NSIG$ but $L^r\not in\NSIG$.  Then $L^r\in\DSIG$.  By
Exercise~5.20, \DSIG\ is closed under reversal, so
$L=(L^r)^r\in\DSIG$, a contradiction.  Hence $L^r\in\NSIG$ whenever
$L\in\NSIG$.

%% 5.23
\item \textbf{\NSIG\ is \emph{not} closed under $P$ (prefix operator).}

Let $K=\{x\in\{\pmb{a},\pmb{b}\}^*\mid |x|_{\pmb{a}}=|x|_{\pmb{b}}\}$.  This
language belongs to \NSIG\ (it is the language used in Lemma~5.4).  But
$P(K)=\{\pmb{a},\pmb{b}\}^*$: given any word $x$, if
$|x|_{\pmb{a}}-|x|_{\pmb{b}}=d$, append $|d|$ copies of the deficient symbol to
obtain a balanced word.  Thus $P(K)$ is FAD, so \NSIG\ is not closed under $P$.

%% 5.24
\item \textbf{\NSIG\ is not closed under $S$ (suffix operator).}

Using the same language $K$, we have $S(K)=\{\pmb{a},\pmb{b}\}^*$: for any word
$x$, prefix it by enough copies of the deficient symbol to balance the counts.
Since $K\in\NSIG$ but $S(K)$ is FAD, \NSIG\ is not closed under $S$.

%% 5.25
\item \textbf{\NSIG\ is not closed under $C$ (center operator).}

Again let $K=\{x\mid |x|_{\pmb{a}}=|x|_{\pmb{b}}\}$.  Then
$C(K)=\{\pmb{a},\pmb{b}\}^*$, because for any word $x$ we can choose
$y=\pmb{a}^m$ and $z=\pmb{b}^{m+|x|_{\pmb{a}}-|x|_{\pmb{b}}}$ with $m$ large
enough so that $yxz$ has equally many $\pmb{a}$s and $\pmb{b}$s.  Hence $C(K)$
is FAD while $K\in\NSIG$.

%% 5.26
\item \textbf{\NSIG\ is not closed under $F$.}

For the same language $K$, every word in $K$ has a proper extension in $K$:
if $x\in K$, then $x\pmb{ab}\in K$.  Therefore $F(K)=\emptyset$, which is FAD.
So \NSIG\ is not closed under $F$.

%% 5.27
\item \textbf{Complementing final states of an NDFA does not give complement language.}

From Exercise~4.8: if $A$ is an NDFA, $L(A^\sim)\neq\sim L(A)$ in general.

\textbf{Example.}  Let $A$ have $\Sigma=\{\pmb{a}\}$, $S=\{s_0,s_1\}$, $S_0=\{s_0\}$, $F=\{s_1\}$,
$\delta(s_0,\pmb{a})=\{s_0,s_1\}$, $\delta(s_1,\pmb{a})=\{s_1\}$.
Then $L(A)=\{\pmb{a}\}^+=\{\pmb{a},\pmb{aa},\ldots\}$.
After complementing: $F^\sim=\{s_0\}$.
$\pmb{a}$: from $s_0$, reach $\{s_0,s_1\}$; $s_0\in F^\sim$, so $\pmb{a}\in L(A^\sim)$.
But $\pmb{a}\in L(A)$, so $\pmb{a}\not in\sim L(A)$.  Thus $\pmb{a}\in L(A^\sim)$ but $\pmb{a}\not in\sim L(A)$, so $L(A^\sim)\neq\sim L(A)$.

%% 5.28
\item \textbf{Complete the proof of Lemma~5.7.}

\begin{proof}
Let
$L=\{x\in\{\pmb{a},\pmb{b}\}^*\mid |x|_{\pmb{a}}\neq |x|_{\pmb{b}}\}$, the
language used in Lemma~5.7.  We show that $L/L=\Sigma^*$.

Take any $x\in\Sigma^*$ and let $d=|x|_{\pmb{a}}-|x|_{\pmb{b}}$.  Define
$y=\pmb{a}^{|d|+1}$.  Then $y\in L$, because it has strictly more $\pmb{a}$s
than $\pmb{b}$s.  Also
\[
|xy|_{\pmb{a}}-|xy|_{\pmb{b}}
= d+(|d|+1)\neq 0,
\]
so $xy\in L$ as well.  Hence $x\in L/L$.  Since $x$ was arbitrary,
$L/L=\Sigma^*$.  Therefore \NSIG\ is not closed under quotient.
\end{proof}

%% 5.29
\item \textbf{Collection closed under union, concatenation, Kleene closure, but not intersection.}

Take $\mathcal{C}=I_\Sigma$, the collection of all infinite languages from
Exercise~5.4.  Exercise~5.4 shows that $I_\Sigma$ is closed under union,
concatenation, and Kleene closure, but not under intersection.  So
$I_\Sigma$ is the required example.

%% 5.30
\item \textbf{Closed under union does not imply closed under intersection.}

\textbf{Counterexample.}  Again use $I_\Sigma$, the infinite languages.
Exercise~5.4(b) shows it is closed under union, while Exercise~5.4(c) shows it
is not closed under intersection.

%% 5.31
\item \textbf{Prove $L(A^\bullet)=L_1\cdot L_2$ (Theorem~5.4 construction).}

\begin{proof}
If $x\in L_1\cdot L_2$, write $x=uv$ with $u\in L_1$ and $v\in L_2$.
Run $A_1$ on $u$ from $s_{0_1}$ to some state of $F_1$.

If $v=\lambda$, then $\lambda\in L_2$, so Theorem~5.4 makes every state of
$F_1$ final in $A^\bullet$; hence the same path accepts $x$.

If $v=\pmb{a}w$ is nonempty, then from the final state reached after $u$ the
definition of $\delta^\bullet$ allows the next input letter $\pmb{a}$ to be
processed exactly as it would be from $s_{0_2}$ in $A_2$.  The rest of $v$ then
follows the transitions of $A_2$, so some path in $A^\bullet$ accepts $x$.
Thus $L_1\cdot L_2\subseteq L(A^\bullet)$.

Conversely, let $x\in L(A^\bullet)$.  An accepting path either stays entirely in
$S_1$ and ends in a state of $F_1$; this can happen only when $\lambda\in L_2$,
so then $x\in L_1=L_1\cdot\lambda\subseteq L_1\cdot L_2$.  Otherwise there is a
first transition that enters $S_2$.  Just before that step, the path is at some
state $f\in F_1$, so the already-read prefix $u$ lies in $L_1$.  The rest of the
input, beginning with the symbol that caused the move into $S_2$, follows
$A_2$ from $s_{0_2}$ to a final state, so it lies in $L_2$.  Hence
$x=uv\in L_1\cdot L_2$.
\end{proof}

%% 5.32
\item \textbf{Prove $L(A_*)=L^*$ (Theorem~5.5 construction).}

\begin{proof}
Since $q_0\in F_*$, the empty word is accepted, so $L^0=\{\lambda\}\subseteq L(A_*)$.

Now let $x=x_1x_2\cdots x_k$ with each $x_i\in L$.  Run the original machine $A$
on $x_1$ from some start state in $S_0$ to a final state in $F$.  Whenever one
factor ends and the next factor begins with a letter $\pmb{a}$, the definition of
$\delta_*$ for final states adds the transitions
$\bigcup_{t\in S_0}\delta(t,\pmb{a})$, so the same input letter can also begin a
new run of $A$ on the next factor.  Repeating this for all factors yields an
accepting path for $x$, so $L^*\subseteq L(A_*)$.

Conversely, consider an accepting path of $A_*$ on some nonempty word $x$.
Because $\delta_*(q_0,\pmb{a})=\emptyset$, the path must begin in some state of
$S_0$.  Each time the path uses one of the extra transitions
$\bigcup_{t\in S_0}\delta(t,\pmb{a})$ from a final state, it starts a fresh copy
of the original machine on the next factor of the input.  Thus the input splits
into finitely many consecutive blocks, each of which labels a path in $A$ from a
start state to a final state; hence each block lies in $L$.  Therefore
$x\in L^*$.
\end{proof}

%% 5.33
\item \textbf{Amplify the proof of Theorem~5.2 (union).}

Theorem~5.2 uses NDFAs
\[
A_1=\langle\Sigma,S_1,S_{0_1},\delta_1,F_1\rangle,\qquad
A_2=\langle\Sigma,S_2,S_{0_2},\delta_2,F_2\rangle,
\]
and the union machine
\[
A^\cup=\langle\Sigma,S_1\cup S_2,S_{0_1}\cup S_{0_2},\delta^\cup,F_1\cup F_2\rangle.
\]
The book's chain of equivalences is:
\begin{align*}
x\in L(A^\cup)
&\Leftrightarrow
(\exists s\in S_{0_1}\cup S_{0_2})
\bigl[\overline{\delta^\cup}(s,x)\cap(F_1\cup F_2)\neq\emptyset\bigr]
\quad\text{(Definition~4.4)} \\
&\Leftrightarrow
(\exists s\in S_{0_1})
\bigl[\overline{\delta^\cup}(s,x)\cap(F_1\cup F_2)\neq\emptyset\bigr] \\
&\hspace{1.5cm}\vee
(\exists s\in S_{0_2})
\bigl[\overline{\delta^\cup}(s,x)\cap(F_1\cup F_2)\neq\emptyset\bigr]
\quad\text{(definition of union)} \\
&\Leftrightarrow
(\exists s\in S_{0_1})
\bigl[\DBAR_1(s,x)\cap(F_1\cup F_2)\neq\emptyset\bigr] \\
&\hspace{1.5cm}\vee
(\exists s\in S_{0_2})
\bigl[\DBAR_2(s,x)\cap(F_1\cup F_2)\neq\emptyset\bigr]
\quad\text{(by the inductive fact about $\overline{\delta^\cup}$)} \\
&\Leftrightarrow
(\exists s\in S_{0_1})
\bigl[\DBAR_1(s,x)\cap F_1\neq\emptyset\bigr] \\
&\hspace{1.5cm}\vee
(\exists s\in S_{0_2})
\bigl[\DBAR_2(s,x)\cap F_2\neq\emptyset\bigr]
\quad\text{(since $S_1\cap S_2=\emptyset$ and $F^\cup=F_1\cup F_2$)} \\
&\Leftrightarrow x\in L(A_1)\vee x\in L(A_2)
\quad\text{(Definition~4.4)} \\
&\Leftrightarrow x\in L_1\cup L_2.
\end{align*}

The missing ingredient in the textbook proof is exactly the inductive claim
\[
\overline{\delta^\cup}(s,x)=
\begin{cases}
\DBAR_1(s,x), & s\in S_1,\\
\DBAR_2(s,x), & s\in S_2,
\end{cases}
\]
which is proved later in Exercise~5.50.

%% 5.34
\item \textbf{Given DFA $A$, define NDFA accepting $L(A)^+$.}

Define $A^+=\langle\Sigma,S,s_0,\delta^+,F\rangle$ where
$\delta^+(s,\pmb{a})=\{\delta(s,\pmb{a})\}\cup\{s_0\mid\delta(s,\pmb{a})\in F\}$.
That is, whenever a transition would lead to a final state, simultaneously allow the machine to restart from $s_0$.  The final states remain $F$.  This accepts $L^+=\bigcup_{k\ge 1}L^k$.

%% 5.35
\item \textbf{Given NDFA $A$, define NDFA accepting $L(A)^+$.}

Same construction as~5.34 adapted to NDFAs:
$A^+=\langle\Sigma,S,S_0,\delta^+,F\rangle$ where for each $s\in S$ and $\pmb{a}\in\Sigma$:
$\delta^+(s,\pmb{a})=\delta(s,\pmb{a})\cup(S_0\text{ if }\delta(s,\pmb{a})\cap F\neq\emptyset)$.
This allows returning to any start state after reaching a final state.

%% 5.36
\item \textbf{If $L$ is FAD, must all subsets of $L$ be FAD?}

\textbf{No.}  Let $\Sigma=\{\pmb{a},\pmb{b}\}$ and $L=\Sigma^*$ (FAD).  The subset
$\{a^nb^n\mid n\ge 0\}$ is not FAD (by the pumping lemma).

%% 5.37
\item \textbf{If $L\in\DSIG$, must all supersets of $L$ be in \DSIG?}

\textbf{No.}  Let $L=\emptyset\in\DSIG$.  The superset $\{a^nb^n\mid n\ge 0\}$ is not FAD.

%% 5.38
\item \textbf{If $L\in\NSIG$, must all subsets be in \NSIG?}

\textbf{No.}  Let
$K=\{x\in\{\pmb{a},\pmb{b}\}^*\mid |x|_{\pmb{a}}=|x|_{\pmb{b}}\}\in\NSIG$.
Then $\{\lambda\}\subseteq K$, but $\{\lambda\}\in\DSIG$.

%% 5.39
\item \textbf{If $L\in\NSIG$, must all supersets be in \NSIG?}

\textbf{No.}  With the same language $K\in\NSIG$, we have
$K\subseteq\{\pmb{a},\pmb{b}\}^*$, but $\{\pmb{a},\pmb{b}\}^*\in\DSIG$.

%% 5.40
\item \textbf{Why is $q_0$ needed in the Kleene closure construction (Theorem~5.5)?}

The extra start state $q_0$ handles the case $\lambda\in L^*$ (which is always
true) independently from the rest of the machine.  Without $q_0$, making $s_0$
both the start and a final state could create spurious acceptance: words that
loop back through the original start states might be accepted even when they are
not valid concatenations of words from $L$.  The new state $q_0$ is needed only
to accept $\lambda$ without changing the behavior of the original start states.

%% 5.41
\item \textbf{Redesign Theorem~5.4 (concatenation) using $\lambda$-transitions.}

Given DFAs $A_1$ and $A_2$, construct $A^\bullet_\lambda=\langle\Sigma,S_1\cup S_2,s^1_0,\delta^\bullet_\lambda,F_2\rangle$ where:
$\delta^\bullet_\lambda(s,\pmb{a})=\delta_1(s,\pmb{a})$ for $s\in S_1$,
$\delta^\bullet_\lambda(s,\pmb{a})=\delta_2(s,\pmb{a})$ for $s\in S_2$, and
$\delta^\bullet_\lambda(f,\lambda)\ni s^2_0$ for each $f\in F_1$.
The $\lambda$-transition from each $f\in F_1$ to $s^2_0$ cleanly separates the two phases.

%% 5.42
\item \textbf{Redesign Theorem~5.5 (Kleene closure) using $\lambda$-transitions, without extra $q_0$.}

Let
\[
A=\langle\Sigma,S,S_0,\delta,F\rangle
\]
be the NDFA for $L$.  For each start state $s\in S_0$, introduce a new copy
$s'$ of that start state.  Let
\[
S'=\{s'\mid s\in S_0\},
\]
with $S'\cap S=\emptyset$.  Define the $\lambda$-NDFA
\[
A^\lambda_*=\langle\Sigma,S\cup S',S',\delta^\lambda_*,F\cup S'\rangle
\]
as follows:
\begin{enumerate}[1.]
  \item For each old state $t\in S$ and each $\pmb{a}\in\Sigma$, keep the old
  transitions:
  \[
  \delta^\lambda_*(t,\pmb{a})=\delta(t,\pmb{a}).
  \]
  \item For each copied start state $s'\in S'$ corresponding to $s\in S_0$, let
  it behave on letters exactly as $s$ did:
  \[
  \delta^\lambda_*(s',\pmb{a})=\delta(s,\pmb{a})
  \qquad (\pmb{a}\in\Sigma).
  \]
  \item For each final state $f\in F$, add $\lambda$-transitions from $f$ to
  every copied start state:
  \[
  \delta^\lambda_*(f,\lambda)\supseteq S'.
  \]
  \item No other $\lambda$-transitions are added.
\end{enumerate}

This makes the copied start states both start states and final states, so
$\lambda$ is accepted immediately, and every time a copy of a word from $L$ is
finished, the machine may jump by a $\lambda$-move to a copied start state and
begin another copy.

\textbf{Why it works.}  If $x=x_1x_2\cdots x_n$ with each $x_i\in L$, start in a
copied start state corresponding to the start used for $x_1$, follow the old
transitions while reading $x_1$, jump by $\lambda$ from the reached final state
to a copied start state for $x_2$, and continue.  Thus every word of $L^*$ is
accepted.

Conversely, every accepting computation starts in a copied start state, and the
only way to enter a copied start state later is by a $\lambda$-move from a final
state.  Therefore the input splits into consecutive blocks, each of which takes
the machine from some copied start state through the old machine to a final
state.  Replacing each copied start state by its original partner in $S_0$
shows that each block lies in $L$.  Hence the whole input lies in $L^*$.

So this redesign does make the single extra start state $q_0$ unnecessary: the
role of ``accept $\lambda$ and restart a new factor'' is carried by copied start
states together with $\lambda$-transitions from final states.

%% 5.43
\item \textbf{Redesign Theorem~5.4 for NDFAs $A_1$ and $A_2$.}

Given NDFAs $A_1=\langle\Sigma,S_1,S^1_0,\delta_1,F_1\rangle$ and $A_2=\langle\Sigma,S_2,S^2_0,\delta_2,F_2\rangle$ with $S_1\cap S_2=\emptyset$:
Define
\[
A^\bullet=\langle\Sigma,S_1\cup S_2,S^1_0,\delta^\bullet,F^\bullet\rangle,
\]
where
\[
F^\bullet=
\left\{
\begin{array}{ll}
F_2, & \lambda\not in L_2,\\
F_1\cup F_2, & \lambda\in L_2,
\end{array}
\right.
\]
and
\[
\delta^\bullet(s,\pmb{a})=
\left\{
\begin{array}{ll}
\delta_1(s,\pmb{a}), & s\in S_1-F_1,\\
\delta_1(s,\pmb{a})\cup\displaystyle\bigcup_{t\in S^2_0}\delta_2(t,\pmb{a}),
  & s\in F_1,\\
\delta_2(s,\pmb{a}), & s\in S_2.
\end{array}
\right.
\]
This is the exact NDFA analogue of Theorem~5.4: when the current state is
already in $F_1$, the same input symbol may also start a run of $A_2$ from any
start state in $S^2_0$.

%% 5.44
\item \textbf{Redesign Theorem~5.5 for DFA $A$.}

If $A=\langle\Sigma,S,s_0,\delta,F\rangle$ is deterministic, the redesigned
machine is the NDFA
\[
A_*=\langle\Sigma,S\cup\{q_0\},\{q_0,s_0\},\delta_*,F\cup\{q_0\}\rangle,
\]
where
\[
\delta_*(s,\pmb{a})=
\left\{
\begin{array}{ll}
\{\delta(s,\pmb{a})\}, & s\in S-F,\\
\{\delta(s,\pmb{a}),\delta(s_0,\pmb{a})\}, & s\in F,\\
\emptyset, & s=q_0.
\end{array}
\right.
\]
This is the DFA-specialized form of Theorem~5.5, with singleton outputs exactly
because the old transition function $\delta$ is deterministic.

%% 5.45
\item \textbf{Comparing $R_L$ and $R_{\sim L}$.}

The right congruence $\equiv_L$ on $\Sigma^*$ is defined by $x\equiv_L y$ iff $\forall z\in\Sigma^*,xz\in L\Leftrightarrow yz\in L$.  The right congruence $\equiv_{\sim L}$ satisfies $x\equiv_{\sim L}y$ iff $\forall z,xz\in\sim L\Leftrightarrow yz\in\sim L$, i.e., $xz\not in L\Leftrightarrow yz\not in L$, i.e., $xz\in L\Leftrightarrow yz\in L$.  So $\equiv_L=\equiv_{\sim L}$: the two right congruences are identical.

This makes sense: if $A$ is a minimal DFA for $L$, then $A^\sim$ (with complemented final states) is a minimal DFA for $\sim L$, and both have the same state set and transition function.  Thus $R_L=R_{\sim L}$.

%% 5.46
\item \begin{enumerate}[(a)]
  \item \textbf{Examples where $R_{L_1\cap L_2}=R_{L_1}\cap R_{L_2}$.}
  Let $L_1=L_2$ be any FAD language.  Then $L_1\cap L_2=L_1$, so $R_{L_1\cap L_2}=R_{L_1}=R_{L_1}\cap R_{L_1}=R_{L_1}\cap R_{L_2}$.

  \item \textbf{Examples where $R_{L_1\cap L_2}\neq R_{L_1}\cap R_{L_2}$.}
  Let $L_1=\{\pmb{a}\}$ and $L_2=\{\pmb{b}\}$ over $\Sigma=\{\pmb{a},\pmb{b}\}$.
  Then $L_1\cap L_2=\emptyset$, so $R_{L_1\cap L_2}=R_\emptyset$ is the
  universal relation.  But $\lambda\not\mathrel{R_{L_1}}\pmb{a}$, because with
  suffix $\lambda$ the word $\pmb{a}$ is in $L_1$ and $\lambda$ is not.  Hence
  \[
  R_{L_1\cap L_2}\neq R_{L_1}\cap R_{L_2}.
  \]
\end{enumerate}

%% 5.47
\item \textbf{\DSIG\ is closed under relative complement.}

\begin{enumerate}[(a)]
  \item \textbf{By theorems.}  $L_1\setminus L_2=L_1\cap\sim L_2$.  Since \DSIG\ is closed under complement (Theorem~5.1) and intersection (Lemma~5.1), $L_1\setminus L_2\in\DSIG$.

  \item \textbf{Machine definition.}  Define $A^-$ as the product of $A_1$ and $A_2^\sim$ (complement of $A_2$), using the intersection machine from Lemma~5.1 with $F^-=F_1\times(S_2\setminus F_2)$.

  \item \textbf{Proof.}  By induction (as in Exercise~5.56/5.57), $\overline{\delta^-}(\langle s^1_0,s^2_0\rangle,x)=\langle\deltabar{1}(s^1_0,x),\deltabar{2}(s^2_0,x)\rangle$.  Then $x\in L(A^-)$ iff $\deltabar{1}(s^1_0,x)\in F_1$ and $\deltabar{2}(s^2_0,x)\not in F_2$ iff $x\in L_1\setminus L_2$. \qed
\end{enumerate}

%% 5.48
\item $\mathfrak{L}_\Sigma$ = languages recognized by $\lambda$-NDFAs.  By Theorem~4.2, $\mathfrak{L}_\Sigma=\DSIG$.  Hence $\mathfrak{L}_\Sigma$ has exactly the same closure properties as \DSIG\ (Theorems~5.1--5.12).

%% 5.49
\item \begin{enumerate}[(a)]
  \item \textbf{Example of $L$ with $\lambda\in L^+$.}  Let $L=\{\lambda\}$.  Then $L^+=\bigcup_{k\ge 1}L^k=\{\lambda\}$, and $\lambda\in L^+$.

  \item \textbf{Three languages with $L^+=L$.}
  \begin{enumerate}[(i)]
    \item $L=\Sigma^*$: $(\Sigma^*)^+=\Sigma^*$.
    \item $L=\{\pmb{a}\}^*$: $(\{\pmb{a}\}^*)^+=\{\pmb{a}\}^*$.
    \item $L=\Sigma^+$: $(\Sigma^+)^+=\Sigma^+$ (since $\Sigma^+\cdot\Sigma^+\subseteq\Sigma^+$).
  \end{enumerate}
\end{enumerate}

%% 5.50
\item \begin{enumerate}[(a)]
  \item \textbf{Prove $\overline{\delta^\cup}(s,x)=\DBAR_i(s,x)$ for $s\in S_i$.}

  \begin{proof}
  By induction on $|x|$.

  \emph{Base ($|x|=0$):}  $\overline{\delta^\cup}(s,\lambda)=\{s\}=\DBAR_i(s,\lambda)$ (by definition of extended transition functions for single start-states, noting NDFAs).

  \emph{Step:}  Assume $\overline{\delta^\cup}(s,x)=\DBAR_i(s,x)$ for $s\in S_i$.  For $x\pmb{a}$:
  \[
    \overline{\delta^\cup}(s,x\pmb{a})
    =\bigcup_{t\in\overline{\delta^\cup}(s,x)}\delta^\cup(t,\pmb{a})
    =\bigcup_{t\in\DBAR_i(s,x)}\delta_i(t,\pmb{a})
    =\DBAR_i(s,x\pmb{a}),
  \]
  where we used the fact that $\DBAR_i(s,x)\subseteq S_i$ (paths from $s\in S_i$ stay in $S_i$, since $\delta_i:S_i\to\wp(S_i)$) and $\delta^\cup(t,\pmb{a})=\delta_i(t,\pmb{a})$ for $t\in S_i$. \qed
  \end{proof}

  \item \textbf{Was this used in Theorem~5.2?}  Yes, implicitly: the proof that $L(A^\cup)=L_1\cup L_2$ relies on the fact that paths starting in $S_1$ stay in $S_1$ and paths starting in $S_2$ stay in $S_2$, which is exactly this lemma.
\end{enumerate}

%% 5.51
\item \begin{enumerate}[(a)]
  \item \textbf{\NSIG\ is not closed under relative complement.}  Let
  $K=\{x\in\{\pmb{a},\pmb{b}\}^*\mid |x|_{\pmb{a}}=|x|_{\pmb{b}}\}\in\NSIG$.
  Then $K-K=\emptyset\in\DSIG$.

  \item \textbf{\NSIG\ is not closed under union.}  If $K\in\NSIG$, then
  $\sim K\in\NSIG$ by Theorem~5.8.  But
  $K\cup\sim K=\Sigma^*\in\DSIG$.

  \item \textbf{\NSIG\ is not closed under concatenation.}  Let
  $E=\{\pmb{a}^n\pmb{b}^n\mid n\ge 1\}$ and define
  $L=\{\pmb{a},\pmb{b}\}^*-E$.  Since $E\in\NSIG$, Theorem~5.8 gives
  $L\in\NSIG$.  Also $\lambda\in L$, and if $x\not in L$ then
  $x=\pmb{a}^n\pmb{b}^n=\pmb{a}\cdot \pmb{a}^{n-1}\pmb{b}^n$ with both factors in
  $L$.  Hence $L\cdot L=\{\pmb{a},\pmb{b}\}^*\in\DSIG$.

  \item \textbf{\NSIG\ is not closed under Kleene closure.}  With the same
  language $L$, we have $\lambda\in L$ and $L\cdot L=\Sigma^*$, so
  $L^*=\Sigma^*\in\DSIG$.

  \item \textbf{If $L\in\NSIG$, then $L^+\in\NSIG$: No.}  For the same language
  $L$, because $\lambda\in L$ we have $L^+=L^*=\Sigma^*$, which is FAD.
\end{enumerate}

%% 5.52
\item \textbf{Why require $S_1\cap S_2=\emptyset$ in Theorem~5.4?}

If $S_1\cap S_2\neq\emptyset$, the transition function $\delta^\bullet$ would be ambiguous: for a state $s\in S_1\cap S_2$, which transitions should apply---those of $A_1$ or $A_2$?  The definition of $\delta^\bullet$ would not be well-defined.  Moreover, the proof that paths through the concatenation machine correctly simulate $A_1$ then $A_2$ relies on the fact that once the computation enters $S_2$, it stays in $S_2$ (and similarly for $S_1$).  If the state sets overlap, a path might inadvertently ``mix'' transitions from $A_1$ and $A_2$.

%% 5.53
\item \textbf{\DSIG\ is closed under $E(L)=\{z\mid\exists y\in\Sigma^+,\exists x\in L\ni z=yx\}$.}

Here $E(L)=\Sigma^+\cdot L$, the set of words obtained by prepending a nonempty
prefix to a word of $L$.

\begin{enumerate}[(a)]
  \item \textbf{Construction.}  Let $B$ be the two-state DFA for $\Sigma^+$:
  start state $p_0$ (nonfinal), final state $p_1$, and for every
  $\pmb{a}\in\Sigma$,
  $\delta_B(p_0,\pmb{a})=p_1$ and $\delta_B(p_1,\pmb{a})=p_1$.  If $A$ accepts
  $L$, apply the concatenation construction of Theorem~5.4 to $B$ and $A$.

  \item \textbf{Proof.}  A word $z$ is accepted by the new machine iff it can be
  written as $z=yx$ where $y$ is accepted by $B$ and $x$ is accepted by $A$.
  That is equivalent to saying $y\in\Sigma^+$ and $x\in L$, so the new machine
  accepts exactly $E(L)$.

  \item Since $\Sigma^+\in\DSIG$ and \DSIG\ is closed under concatenation,
  $E(L)=\Sigma^+\cdot L\in\DSIG$. \qed
\end{enumerate}

%% 5.54
\item \textbf{\DSIG\ is closed under $B(L)=\{z\mid\exists x\in L,\exists y\in\Sigma^*\ni z=xy\}$.}

$B(L)$ is the set of all words that have a prefix in $L$.

\begin{enumerate}[(a)]
  \item \textbf{Construction.}  Given
  $A=\langle\Sigma,S,s_0,\delta,F\rangle$ for $L$, build
  \[
  A'=\langle\Sigma,S\times\{0,1\},\langle s_0,0\rangle,\delta',S\times\{1\}\rangle,
  \]
  where the second component records whether some prefix read so far belongs to
  $L$.  Define
  \[
  \delta'(\langle s,b\rangle,a)=\langle t,b'\rangle,
  \]
  where $t=\delta(s,a)$ and
  \[
  b'=
  \begin{cases}
  1 & \text{if } b=1 \text{ or } t\in F,\\
  0 & \text{otherwise.}
  \end{cases}
  \]
  Thus once the flag becomes $1$, it stays $1$ forever.

  \item \textbf{Proof:}  By induction on the length of the input, after reading a
  word $z$ the machine $A'$ is in state
  \[
  \langle \deltabar{}(s_0,z), 1\rangle
  \]
  exactly when some prefix of $z$ is in $L$, and otherwise in
  \[
  \langle \deltabar{}(s_0,z), 0\rangle.
  \]
  Therefore $A'$ accepts exactly those words having a prefix in $L$, that is,
  exactly the words in $B(L)$.

  \item \qed
\end{enumerate}

%% 5.55
\item \textbf{\DSIG\ is closed under $M(L)=\{z\mid\exists x\in L,\exists y\in\Sigma^+\ni z=xy\}$.}

$M(L)=\{z\mid z=xy,x\in L,y\in\Sigma^+\}=L\cdot\Sigma^+$.  By closure under concatenation, $L\cdot\Sigma^+\in\DSIG$.

\begin{enumerate}[(a)]
  \item \textbf{Construction:} Concatenate $A$ (for $L$) with a DFA for $\Sigma^+$ (which is a two-state machine: start nonfinal, self-loop on $\Sigma$ to final, final self-loops on $\Sigma$).
  \item \textbf{Proof:} By Theorem~5.4, $M(L)=L(A^\bullet)\in\DSIG$.
  \item \qed
\end{enumerate}

%% 5.56
\item \textbf{Prove $\overline{\delta^\cap}(\langle s,t\rangle,x)=\langle\DBAR_1(s,x),\DBAR_2(t,x)\rangle$.}

\begin{proof}
By induction on $|x|$.

\emph{Base:}  $\overline{\delta^\cap}(\langle s,t\rangle,\lambda)=\langle s,t\rangle=\langle\DBAR_1(s,\lambda),\DBAR_2(t,\lambda)\rangle$.

\emph{Step:}  Assume the result for $x$.  Then:
\begin{align*}
  \overline{\delta^\cap}(\langle s,t\rangle,x\pmb{a})
  &=\delta^\cap(\overline{\delta^\cap}(\langle s,t\rangle,x),\pmb{a}) \\
  &=\delta^\cap(\langle\DBAR_1(s,x),\DBAR_2(t,x)\rangle,\pmb{a}) \\
  &=\langle\delta_1(\DBAR_1(s,x),\pmb{a}),\delta_2(\DBAR_2(t,x),\pmb{a})\rangle \\
  &=\langle\DBAR_1(s,x\pmb{a}),\DBAR_2(t,x\pmb{a})\rangle. \qed
\end{align*}
\end{proof}

%% 5.57
\item \textbf{Prove $L(A^\cap)=L_1\cap L_2$.}

\begin{proof}
By Exercise~5.56:
\begin{align*}
  x\in L(A^\cap)
  &\Leftrightarrow\overline{\delta^\cap}(\langle s^1_0,s^2_0\rangle,x)\in F^\cap \\
  &\Leftrightarrow\langle\DBAR_1(s^1_0,x),\DBAR_2(s^2_0,x)\rangle\in F_1\times F_2 \\
  &\Leftrightarrow\DBAR_1(s^1_0,x)\in F_1\wedge\DBAR_2(s^2_0,x)\in F_2 \\
  &\Leftrightarrow x\in L_1\wedge x\in L_2 \\
  &\Leftrightarrow x\in L_1\cap L_2. \qed
\end{align*}
\end{proof}

%% 5.58
\item \textbf{Prove Theorem~5.6 (\DSIG\ closed under $Z$).}

\begin{proof}
Let $A=\langle\Sigma,S,s_0,\delta,F\rangle$ be a DFA for $L$.  Build a
$\lambda$-NDFA $A_Z$ by adding, for every ordinary transition
$s\xrightarrow{\pmb{a}}t$ of $A$, a matching $\lambda$-transition
$s\xrightarrow{\lambda}t$.  Then a path in $A_Z$ may either read a letter of the
input or skip a letter that was present in some word of $L$.

Thus $x$ is accepted by $A_Z$ iff there exists some $w\in L$ from which $x$ can
be obtained by deleting zero or more letters, i.e., iff $x\in Z(L)$.  Hence
$Z(L)\in\DSIG$.
\end{proof}

%% 5.59
\item \textbf{Prove Theorem~5.7 (\DSIG\ closed under $Y$).}

\begin{proof}
Let $A=\langle\Sigma,S,s_0,\delta,F\rangle$ be a DFA for $L$.  Following the text
after Theorem~5.7, build two copies of $A$.  Being in the first copy means that
no letter has yet been deleted; being in the second copy means that exactly one
letter has been deleted.  Ordinary transitions stay within each copy, and for
every transition $s\xrightarrow{\pmb{a}}t$ in the first copy add a
$\lambda$-transition from the first-copy state $s$ to the second-copy state $t$.

An accepting path therefore behaves like a path of $A$ on some word of $L$, with
exactly one letter skipped at the moment when the path moves from the first copy
to the second.  So the resulting machine accepts exactly $Y(L)$, and therefore
$Y(L)\in\DSIG$.
\end{proof}

%% 5.60
\item \begin{enumerate}[(a)]
  \item \textbf{Alternative construction for Theorem~5.7 without $\lambda$-moves.}

  Let $A=\langle\Sigma,S,s_0,\delta,F\rangle$ be a DFA for $L$.  Use state set
  $S\times\{0,1\}$, where the second component records whether the one deletion
  has already occurred.  The start state is $\langle s_0,0\rangle$.  Final states
  are
  \[
  \{\langle s,1\rangle\mid s\in F\}\cup
  \{\langle s,0\rangle\mid(\exists \pmb{a}\in\Sigma)\ \delta(s,\pmb{a})\in F\}.
  \]
  On input $\pmb{a}$, define
  \[
  \delta'(\langle s,0\rangle,\pmb{a})=
  \{\langle \delta(s,\pmb{a}),0\rangle\}\cup
  \{\langle \delta(\delta(s,\pmb{b}),\pmb{a}),1\rangle\mid \pmb{b}\in\Sigma\},
  \]
  and
  \[
  \delta'(\langle s,1\rangle,\pmb{a})=
  \{\langle \delta(s,\pmb{a}),1\rangle\}.
  \]

  \item \textbf{Proof:}  In the $0$-copy the machine either reads the current
  input symbol normally, or guesses that one letter $\pmb{b}$ of the original
  word was deleted immediately before it and moves into the $1$-copy.  Once in
  the $1$-copy, the rest of the input is processed exactly as in $A$.  The extra
  final states in the $0$-copy handle the case where the deleted letter was the
  last letter of the original word.  Hence this NDFA accepts exactly $Y(L)$.
\end{enumerate}

%% 5.61
\item \textbf{\DSIG\ is closed under $W(L)$ (delete one or more letters).}

\begin{enumerate}[(a)]
  \item \textbf{Construction.}  Given DFA $A$ for $L$, build $\lambda$-NDFA $A'$ by adding, for each transition $s\xrightarrow{\pmb{a}}t$, an optional $\lambda$-transition $s\xrightarrow{\lambda}t$ (allowing deletion of $\pmb{a}$).  Require at least one deletion by adding a flag bit.

  More precisely: $A'$ has state set $S\times\{0,1\}$ where the second
  component tracks whether at least one letter has been deleted.  Initial
  state: $\langle s_0,0\rangle$.  Final states:
  $\{\langle f,1\rangle\mid f\in F\}$.  For each $c\in\Sigma$,
  \[
  \delta'(\langle s,b\rangle,c)=\{\langle\delta(s,c),b\rangle\},
  \qquad
  \delta'(\langle s,b\rangle,\lambda)
    =\{\langle\delta(s,c),1\rangle\mid c\in\Sigma\}.
  \]
  The ordinary transition reads the next input symbol, while the
  $\lambda$-transition simulates deleting one symbol $c$ from the original word
  and moves into the ``already deleted'' copy.

  \item \textbf{Proof.}  A word $z$ is in $W(L)$ iff $z$ can be obtained by deleting at least one letter from some $w\in L$.  A path in $A'$ simulates a path in $A$ but uses $\lambda$-transitions to skip letters; the flag ensures at least one skip. Hence $L(A')=W(L)$.

  \item Since $A'$ is a $\lambda$-NDFA, $W(L)\in\DSIG$. \qed
\end{enumerate}

%% 5.62
\item \textbf{\DSIG\ is closed under $V(L)$ (delete odd-positioned letters).}

\begin{enumerate}[(a)]
  \item \textbf{Construction.}  Let
  $A=\langle\Sigma,S,s_0,\delta,F\rangle$.  Build the $\lambda$-NDFA
  \[
  A'=\langle\Sigma,S\times\{0,1\},\langle s_0,0\rangle,\delta',F\times\{0,1\}\rangle,
  \]
  where the second component records the parity of the next position in the
  original word: $0$ means odd, $1$ means even.  For each $s\in S$ and
  $a\in\Sigma$,
  \[
  \delta'(\langle s,1\rangle,a)=\{\langle\delta(s,a),0\rangle\},
  \]
  because even-positioned letters are kept, and
  \[
  \delta'(\langle s,0\rangle,\lambda)
  =\{\langle\delta(s,c),1\rangle\mid c\in\Sigma\},
  \]
  because an odd-positioned letter is deleted and may be any symbol $c$.

  \item \textbf{Proof.}  A computation of $A'$ alternates between deleting one
  symbol of the original word at every odd position and reading one input symbol
  at every even position.  Thus the visible input is exactly the word obtained
  by deleting the odd-positioned letters of some word accepted by $A$.
  Conversely, if a word $z$ arises from deleting the odd-positioned letters of
  some $w\in L$, then $A'$ can guess those deleted odd-positioned letters by
  $\lambda$-moves and read precisely the surviving even-positioned letters.  So
  $L(A')=V(L)$.

  \item $L(A')\in\DSIG$. \qed
\end{enumerate}

%% 5.63
\item \textbf{\DSIG\ is closed under $U(L)$ (delete even-positioned letters).}

Let
\[
A=\langle\Sigma,S,s_0,\delta,F\rangle.
\]
Build the $\lambda$-NDFA
\[
A'=\langle\Sigma,S\times\{0,1\},\langle s_0,0\rangle,\delta',F\times\{0,1\}\rangle,
\]
where again $0$ means the next position in the original word is odd and $1$
means even.  This time odd-positioned letters are kept and even-positioned
letters are deleted:
\[
\delta'(\langle s,0\rangle,a)=\{\langle\delta(s,a),1\rangle\}
\qquad (a\in\Sigma),
\]
\[
\delta'(\langle s,1\rangle,\lambda)
=\{\langle\delta(s,c),0\rangle\mid c\in\Sigma\}.
\]
Exactly the same reasoning as in 5.62 shows that $A'$ accepts precisely those
words obtained from words of $L$ by deleting the even-positioned letters.
Hence $U(L)\in\DSIG$. \qed

%% 5.64
\item \textbf{\DSIG\ is closed under $T(L)$ (delete every 3rd, 6th, ... letter).}

\begin{enumerate}[(a)]
  \item \textbf{Construction.}  Let
  $A=\langle\Sigma,S,s_0,\delta,F\rangle$.  Build the $\lambda$-NDFA
  \[
  A'=\langle\Sigma,S\times\{1,2,0\},\langle s_0,1\rangle,\delta',F\times\{1,2,0\}\rangle,
  \]
  where the second component records the next position modulo $3$.  For each
  $s\in S$ and $a\in\Sigma$,
  \[
  \delta'(\langle s,1\rangle,a)=\{\langle\delta(s,a),2\rangle\},
  \qquad
  \delta'(\langle s,2\rangle,a)=\{\langle\delta(s,a),0\rangle\},
  \]
  because positions $1$ and $2$ modulo $3$ are kept, while
  \[
  \delta'(\langle s,0\rangle,\lambda)
  =\{\langle\delta(s,c),1\rangle\mid c\in\Sigma\},
  \]
  because every third letter is deleted and may be any symbol $c$.

  \item \textbf{Proof.}  The machine simulates $A$ on an original word while the
  counter records which positions are first, second, and third in each block of
  three.  At positions congruent to $1$ or $2$ modulo $3$, it reads the next
  input symbol and processes it through $\delta$; at positions congruent to $0$
  modulo $3$, it guesses and deletes one original symbol by a $\lambda$-move.
  Thus the visible input is exactly the word obtained from the original one by
  deleting every $3$rd, $6$th, $9$th, \dots letter.  Hence $L(A')=T(L)$.

  \item $L(A')\in\DSIG$. \qed
\end{enumerate}

%% 5.65
\item \textbf{$I(L)=L\cap P$ where $P=\{x\mid|x|\text{ is prime}\}$.}

\begin{enumerate}[(a)]
  \item \textbf{\DSIG\ not closed under $I$.}  Let $\Sigma=\{\pmb{a}\}$ and
  $L=\Sigma^*$.  Then
  \[
  I(L)=P=\{\pmb{a}^q\mid q\text{ prime}\}.
  \]
  We show $P\not in\DSIG$ by the pumping lemma.  Assume $P$ were regular, and
  let $p$ be a pumping length.  Choose a prime $q>p$, and consider
  $\pmb{a}^q\in P$.  Write
  \[
  \pmb{a}^q=uvw
  \]
  with $|uv|\le p$ and $|v|\ge 1$.  Since the word is unary, let $|v|=m\ge 1$.
  Pump with $i=q+1$.  Then
  \[
  |uv^{q+1}w|=q+q m=q(1+m).
  \]
  This number is composite, because $q>1$ and $1+m>1$.  Hence
  $uv^{q+1}w\not in P$, contradicting the pumping lemma.  Therefore
  $P\not in\DSIG$.

  \item \textbf{$F_\Sigma$ closed under $I$.}  $I(L)=L\cap P\subseteq L$; if $L$ is finite, so is $I(L)$.

  \item \textbf{$C_\Sigma$ closed under $I$? No.}  $L=\Sigma^*\setminus F$ for finite $F$.  $I(L)=P\setminus F'$ for some finite $F'$.  If $P$ is cofinite: $P$ is not cofinite (infinitely many non-primes), so $P\setminus F'$ is not cofinite.  \textbf{Not closed.}

  \item \textbf{$B_\Sigma$ closed under $I$? No.}  $I(\Sigma^*)=P\not in B_\Sigma$ (not finite, not cofinite).

  \item \textbf{$I_\Sigma$ closed under $I$? No.}  Let $\Sigma=\{\pmb{a}\}$ and
  $L=\{\pmb{a}^{2n}\mid n\ge 1\}$.  Then $L\in I_\Sigma$, but
  \[
  I(L)=L\cap P=\{\pmb{aa}\},
  \]
  since $2$ is the only even prime.  Thus $I(L)\not in I_\Sigma$.

  \item \textbf{$J_\Sigma$ closed under $I$?}  $I(L)=L\cap P$.  If $\sim L$ is infinite, $\sim(L\cap P)\supseteq\sim P$ which is infinite (non-primes).  \textbf{Yes, closed.}

  \item \textbf{$\pmb{E}$ closed under $I$?}  If $\pmb{abba}\in L$ and $|\pmb{abba}|=4$ is not prime, then $\pmb{abba}\not in I(L)$.  So $I(L)$ might not contain $\pmb{abba}$.  \textbf{Not closed.}

  \item \textbf{\NSIG\ closed under $I$? No.}  Over $\Sigma=\{\pmb{a}\}$, let
  $P=\{\pmb{a}^p\mid p\text{ prime}\}$ and
  $L=\{\pmb{a}\}^*-P$.  Since $P\not in\DSIG$, Theorem~5.8 implies
  $L\in\NSIG$.  But
  \[
  I(L)=L\cap P=\emptyset\in\DSIG.
  \]
  Hence \NSIG\ is not closed under $I$.
\end{enumerate}

%% 5.66
\item $C$: all languages over $\{\pmb{a},\pmb{b}\}$ not containing $\lambda$.

\begin{enumerate}[(a)]
  \item \textbf{Complementation: No.}  $\sim L$ contains $\lambda$ (since $\lambda\not in L$ but $\lambda\in\Sigma^*$), so $\sim L\not in C$.

  \item \textbf{Union: Yes.}  If $\lambda\not in L_1$ and $\lambda\not in L_2$, then $\lambda\not in L_1\cup L_2$.

  \item \textbf{Intersection: Yes.}  $\lambda\not in L_1\cap L_2\subseteq L_1$.

  \item \textbf{Concatenation: Yes.}  If $\lambda\not in L_1$ and $\lambda\not in L_2$: every word in $L_1\cdot L_2$ is of the form $xy$ with $x\in L_1,y\in L_2$, and $x\neq\lambda,y\neq\lambda$, so $xy\neq\lambda$.

  \item \textbf{Kleene closure: No.}  $\lambda\in L^*$ always.

  \item \textbf{Relative complement: Yes.}  $\lambda\not in L_1\setminus L_2\subseteq L_1$.

  \item \textbf{$L^+$: Yes.}  $L^+=\bigcup_{k\ge 1}L^k$.  Each $L^k$ contains words of positive length (since $\lambda\not in L$), so $\lambda\not in L^+$.
\end{enumerate}

%% 5.67
\item \begin{enumerate}[(a)]
  \item \textbf{\DSIG\ closed under finite union.}
  \begin{enumerate}[(i)]
    \item \emph{By theorems and induction:}  Base: $\{L\}\subseteq\DSIG$ trivially.  Step: if $L_1\cup L_2\cup\cdots\cup L_k\in\DSIG$ and $L_{k+1}\in\DSIG$, then $(L_1\cup\cdots\cup L_k)\cup L_{k+1}\in\DSIG$ by Theorem~5.2.
    \item \emph{By construction:}  Build $A^{(k)}$ iteratively by the union construction of Exercise~5.14; each step is deterministic.
  \end{enumerate}
  \item \textbf{\DSIG\ is NOT closed under infinite union.}  Each singleton $\{a^nb^n\}\in\DSIG$ (finite languages are regular), but $\bigcup_{n\ge 0}\{a^nb^n\}=\{a^nb^n\mid n\ge 0\}\not in\DSIG$ (not regular, by pumping).
\end{enumerate}

%% 5.68
\item \begin{enumerate}[(a)]
  \item \textbf{Three homomorphisms under which \NSIG\ is not closed} (over $\Sigma=\{\pmb{a},\pmb{b}\}$):
  Use the non-FAD language
  $K=\{x\in\{\pmb{a},\pmb{b}\}^*\mid |x|_{\pmb{a}}=|x|_{\pmb{b}}\}$ from
  Lemma~5.4.
  \begin{enumerate}[(i)]
    \item $\xi_1(\pmb{a})=\pmb{a},\ \xi_1(\pmb{b})=\pmb{a}$.  By Lemma~5.4,
    $\overline{\xi_1}(K)$ is the regular language of even-length $\pmb{a}$-strings.
    \item $\xi_2(\pmb{a})=\lambda,\ \xi_2(\pmb{b})=\lambda$.  Then
    $\overline{\xi_2}(K)=\{\lambda\}\in\DSIG$.
    \item $\xi_3(\pmb{a})=\pmb{a},\ \xi_3(\pmb{b})=\lambda$.  For
    $L=\{\pmb{a}^n\pmb{b}^n\mid n\ge 0\}\in\NSIG$, we get
    $\overline{\xi_3}(L)=\{\pmb{a}^n\mid n\ge 0\}\in\DSIG$.
  \end{enumerate}

  \item \textbf{Three homomorphisms under which \NSIG\ is closed:}
  \begin{enumerate}[(i)]
    \item The identity map: $\psi_1(\pmb{a})=\pmb{a},\ \psi_1(\pmb{b})=\pmb{b}$.
    \item The swap map: $\psi_2(\pmb{a})=\pmb{b},\ \psi_2(\pmb{b})=\pmb{a}$.
    \item The doubling map: $\psi_3(\pmb{a})=\pmb{aa},\ \psi_3(\pmb{b})=\pmb{bb}$.
  \end{enumerate}
\noindent
For $\psi_1$ and $\psi_2$, if $\overline{\psi_i}(L)$ were FAD then applying the
inverse letter permutation would make $L$ FAD.  For $\psi_3$, the map is
injective, so
\[
L=\psi_3^{-1}(\overline{\psi_3}(L)).
\]
Thus, if $\overline{\psi_3}(L)$ were FAD, then Theorem~5.10 (closure under
inverse homomorphism) would imply that $L$ is FAD as well.
\end{enumerate}

%% 5.69
\item Over $\Sigma=\{\pmb{a}\}$ there is only one homomorphism under which
\NSIG\ is not closed, namely the erasing homomorphism
$\psi(\pmb{a})=\lambda$.  For any non-FAD unary language $L$,
$\overline{\psi}(L)=\{\lambda\}\in\DSIG$.

For every other unary homomorphism, $\psi(\pmb{a})=\pmb{a}^k$ with $k\ge 1$.
Such a homomorphism is injective, so if $\overline{\psi}(L)$ were FAD, then
\[
L=\overline{\psi^{-1}}(\overline{\psi}(L))
\]
would be FAD by Theorem~5.10.  Therefore \NSIG\ is closed under every nonerasing
unary homomorphism.  So the answer is \textbf{no}: there are not two different
unary homomorphisms under which \NSIG\ is not closed.

%% 5.70
\item \begin{enumerate}[(a)]
  \item \textbf{Prove $\overline{\delta^\psi}(s,x)=\DBAR(s,\PBAR(x))$.}

  \begin{proof}
  By induction on $|x|$.

  \emph{Base:} $\overline{\delta^\psi}(s,\lambda)=s=\DBAR(s,\lambda)=\DBAR(s,\PBAR(\lambda))$ (since $\psi(\lambda)=\lambda$).

  \emph{Step:} $\overline{\delta^\psi}(s,x\pmb{b})=\delta^\psi(\overline{\delta^\psi}(s,x),\pmb{b})=\deltabar{}(\DBAR(s,\PBAR(x)),\psi(\pmb{b}))=\DBAR(s,\PBAR(x)\psi(\pmb{b}))=\DBAR(s,\PBAR(x\pmb{b}))$. \qed
  \end{proof}

  \item \textbf{Proof of Theorem~5.10:} $y\in L(A^\psi)$ iff $\overline{\delta^\psi}(s_0,y)\in F$ iff $\DBAR(s_0,\psi(y))\in F$ iff $\psi(y)\in L$ iff $y\in\psi^{-1}(L)$. \qed
\end{enumerate}

%% 5.71
\item \begin{enumerate}[(a)]
  \item \textbf{Apparent contradiction resolved.}  $\xi(L)$ is the set of even-length $\pmb{a}$-strings; $\xi^{-1}(\xi(L))$ is the set of all strings $x$ with $\xi(x)\in\xi(L)$, not $L$ itself.  In general, $\xi^{-1}(\xi(L))\supseteq L$ and can be much larger; it is not necessarily $L$.  The inverse image lands in $\DSIG$ by Theorem~5.10, but it's $\xi^{-1}(\xi(L))$, not $L$.

  \item \textbf{Example where $\PBAR(\psi^{-1}(L))\neq L$:} Take $\psi(\pmb{a})=\pmb{a}\pmb{a}$, $L=\{\pmb{a}\}$ (singleton).  Then $\psi^{-1}(L)=\emptyset$ (no word maps to $\pmb{a}$ under $\psi$, since $\psi$ always produces even-length strings), and $\PBAR(\emptyset)=\emptyset\neq L$.

  \item \textbf{Example where $\psi^{-1}(\PBAR(L))\neq L$:} Same: $\psi^{-1}(\PBAR(\{\pmb{a}\}))=\psi^{-1}(\{\pmb{a}\})=\emptyset\neq\{\pmb{a}\}=L$.

  \item \textbf{Prove $\PBAR(\psi^{-1}(L))\subseteq L$:} Let $y\in\PBAR(\psi^{-1}(L))$.  Then $y=\psi(x)$ for some $x\in\psi^{-1}(L)$.  Since $x\in\psi^{-1}(L)$, $\psi(x)\in L$, i.e., $y\in L$. \qed

  \item \textbf{Prove $L\subseteq\psi^{-1}(\PBAR(L))$:} Let $x\in L$.  Then $\psi(x)\in\PBAR(L)$ (since $\psi(x)$ is the image of some element of $L$, namely $x$ itself).  Hence $x\in\psi^{-1}(\PBAR(L))$. \qed
\end{enumerate}

%% 5.72
\item \textbf{\DSIG\ is closed under $L^e$ (last halves of even-length words in $L$).}

$L^e=\{x\mid\exists y\in\Sigma^*\ni yx\in L\wedge|x|=|y|\}$.

\begin{proof}
Observe that
\[
x\in L^e
\iff \exists y\in\Sigma^*\ (yx\in L\ \wedge\ |y|=|x|)
\iff x^r\in (L^r)^b.
\]
Indeed, $yx\in L$ with $|y|=|x|$ iff $x^ry^r\in L^r$ with
$|x^r|=|y^r|$, which says exactly that $x^r$ is a first half of some word in
$L^r$.

Therefore
\[
L^e=\bigl((L^r)^b\bigr)^r.
\]
Exercise~5.20 shows that \DSIG\ is closed under reversal, and
Theorem~5.11 shows that it is closed under $b$.  Hence \DSIG\ is closed under
$e$.
\end{proof}

%% 5.73
\item \textbf{Show that redefining final states to merely ``midway'' states is incorrect.}

Example~5.25 gives a concrete counterexample.  For the highlighted state $q$,
\[
I(A,q)=\{(\pmb{abc})^n\mid n\ge 1\}
\qquad\text{and}\qquad
T(A,q)=\{\pmb{a}^{2n}\mid n\ge 1\}.
\]
So a word can use $q$ as a midpoint only when its first half has a length that
occurs in both sets, namely a length divisible by $6$.  The proof of
Theorem~5.11 therefore keeps only
\[
I(A,q)\cap\overline{\psi^{-1}}(\PBAR(I(A,q))\cap\PBAR(T(A,q)))
=\{(\pmb{abc})^{2n}\mid n\ge 1\}.
\]
If we simply declared $q$ to be a final state, the resulting automaton would
accept all of $I(A,q)$, including $\pmb{abc}$.  But $\pmb{abc}$ is not a first
half of any word in $L$, because after reaching $q$ there is no continuation of
length $3$ to a final state.  Hence ``make the midway states final'' accepts too
many words and is incorrect.

%% 5.74
\item \textbf{Flaw in the argument that $K$ is not FAD because $M$ is not FAD.}

Theorem~5.9 states that \DSIG\ is closed under homomorphism: if $L$ is FAD and
$\psi$ is a homomorphism, then $\psi(L)$ is FAD.  This does \emph{not} mean that
if $L$ is not FAD, then $\psi(L)$ is not FAD.  A homomorphism can collapse the
structure that made $M$ non-FAD, yielding a simpler FAD image $K=\psi(M)$.
Hence the fact that $M\not in\DSIG$ says nothing by itself about whether
$K=\psi(M)$ is FAD.

%% 5.75
\item \textbf{Prove Theorem~5.9 (\DSIG\ closed under homomorphism).}

\begin{proof}
Let $A=\langle\Sigma,S,s_0,\delta,F\rangle$ be a DFA for $L$, and let
$\psi:\Sigma\to\Sigma^*$ be a homomorphism.  For each transition
$s\xrightarrow{\pmb{a}}t$ of $A$, replace that edge by a path labeled
$\psi(\pmb{a})$.  If $\psi(\pmb{a})=\lambda$, use a $\lambda$-transition from
$s$ to $t$; if $\psi(\pmb{a})=\pmb{b}_1\cdots\pmb{b}_k$ with $k\ge 1$, insert
$k-1$ new states so that the new path reads
$\pmb{b}_1,\pmb{b}_2,\dots,\pmb{b}_k$.

The resulting finite automaton accepts exactly the homomorphic images of words in
$L$, namely $\overline{\psi}(L)$.  Hence $\overline{\psi}(L)\in\DSIG$, proving
Theorem~5.9.
\end{proof}

%% 5.76
\item \textbf{Homomorphism capitalizing lowercase ASCII letters.}

Define $\psi:\text{ASCII}^*\to\text{ASCII}^*$ by
$\psi(c)=\begin{cases}\text{uppercase}(c) & \text{if }c\in\{a,\ldots,z\} \\ c & \text{otherwise}\end{cases}$
(extended letter-by-letter to words).  This is a homomorphism since $\psi(xy)=\psi(x)\psi(y)$ follows from the letter-by-letter definition.

%% 5.77
\item \begin{enumerate}[(a)]
  \item \textbf{Prove $L(A^{^/})=L_1/L_2$.}

  $L_1/L_2=\{x\mid\exists y\in L_2\ni xy\in L_1\}$.
  $F^{^/}=\{t\in S_1\mid\exists y\in\Sigma^*\ni y\in L_2\wedge\DBAR_1(t,y)\in F_1\}$.

  \begin{proof}
  $x\in L(A^{^/})$ iff $\DBAR_1(s_{0_1},x)\in F^{^/}$
  iff $\exists y\in L_2\ni\DBAR_1(\DBAR_1(s_{0_1},x),y)\in F_1$
  iff $\exists y\in L_2\ni\DBAR_1(s_{0_1},xy)\in F_1$
  iff $\exists y\in L_2\ni xy\in L_1$
  iff $x\in L_1/L_2$. \qed
  \end{proof}

  \item \textbf{Algorithm to compute $F^{^/}$.}

  For each state $t\in S_1$, determine whether $\exists y\in L_2\ni\DBAR_1(t,y)\in F_1$:
  \begin{enumerate}[1.]
    \item Build the product machine $A_1\times A_2$: state set $S_1\times S_2$, start at $\langle t,s_{0_2}\rangle$ for each $t$.
    \item In this product, run $A_2$ on the input and simultaneously advance $A_1$.
    \item State $t\in F^{^/}$ iff some state $\langle f_1,f_2\rangle$ with $f_1\in F_1$ and $f_2\in F_2$ is reachable from $\langle t,s_{0_2}\rangle$ in the product.
  \end{enumerate}
  This is a reachability computation in a finite graph, computable in $O(|S_1|\cdot|S_2|)$ time.
\end{enumerate}

%% 5.78
\item \textbf{Extend DFA over $\Sigma_1$ to handle $\Sigma_1\cup\Sigma_2$.}

\begin{enumerate}[(a)]
  \item Define $A'=\langle\Sigma_1\cup\Sigma_2,S\cup\{d\},s_0,\delta',F\rangle$ where $d$ is a new dead state and
  $\delta'(s,\pmb{a})=\begin{cases}\delta(s,\pmb{a}) & \pmb{a}\in\Sigma_1 \\ d & \pmb{a}\in\Sigma_2\setminus\Sigma_1\end{cases}$
  and $\delta'(d,\pmb{a})=d$ for all $\pmb{a}$.  Final states: $F$ (unchanged).

  \item \textbf{Proof.}  If $w\in(\Sigma_1\cup\Sigma_2)^*$ and $w$ contains any letter from $\Sigma_2\setminus\Sigma_1$, then $A'$ goes to $d$ and stays there, rejecting $w$.  If $w\in\Sigma_1^*$, then $A'$ behaves identically to $A$.  Hence $L(A')=L(A)$ (viewed as a subset of $(\Sigma_1\cup\Sigma_2)^*$). \qed
\end{enumerate}

%% 5.79
\item \textbf{If $\mathcal{S}$ is closed under union, concatenation, Kleene closure, and is infinite, must every language in $\mathcal{S}$ be FAD?}

\textbf{No.}  The collection $I_\Sigma$ of all infinite languages is infinite,
and Exercise~5.4 shows that it is closed under union, concatenation, and Kleene
closure.  But $I_\Sigma$ contains languages in \NSIG, so not every language in
such a collection must be FAD.

%% 5.80
\item \textbf{If $\mathcal{S}$ is closed under union, concatenation, Kleene closure, and is finite, must every language be FAD?}

\textbf{No.}  Let
\[
K=\{x\in\{\pmb{a},\pmb{b}\}^*\mid |x|_{\pmb{a}}=|x|_{\pmb{b}}\},
\]
which is in \NSIG\ by Lemma~5.4, and define
\[
\mathcal{S}=\{\emptyset,\{\lambda\},K\}.
\]
This collection is finite.  It is closed under union because
$\{\lambda\}\subseteq K$, so every union is still one of the three listed
languages.  It is closed under concatenation because
\[
\emptyset\cdot L=\emptyset,\qquad \{\lambda\}\cdot L=L\cdot\{\lambda\}=L,
\qquad K\cdot K=K
\]
(the numbers of $\pmb{a}$s and $\pmb{b}$s add, and $\lambda\in K$).  It is also
closed under Kleene closure:
\[
\emptyset^*=\{\lambda\},\qquad \{\lambda\}^*=\{\lambda\},\qquad K^*=K.
\]
So a finite collection with these closure properties need not consist entirely of
FAD languages.

%% 5.81 (bonus)
\item \textbf{If $u\circ u=\text{id}$ and \DSIG\ is closed under $u$, then \NSIG\ is closed under $u$.}

\begin{proof}
Let $L\in\NSIG$.  Suppose, toward a contradiction, that $u(L)\in\DSIG$.
Because \DSIG\ is closed under $u$, applying $u$ again gives
\[
u(u(L))\in\DSIG.
\]
But $u\circ u=\mathrm{id}$, so $u(u(L))=L$.  Hence $L\in\DSIG$, contradicting
$L\in\NSIG$.  Therefore $u(L)\in\NSIG$ whenever $L\in\NSIG$, so \NSIG\ is closed
under $u$.
\end{proof}

\end{enumerate}
